% Standalone problem document, generated from the adjacent README.md.
% Compile with XeLaTeX. Mathematical content is maintained in Markdown.
\documentclass[11pt,a4paper]{article}
\usepackage[fontsize=10.5pt]{scrextend}
\usepackage[margin=22mm,headheight=15pt,headsep=9mm,footskip=11mm]{geometry}
\usepackage{fontspec}
\setmainfont{texgyrepagella}[Extension=.otf,UprightFont=*-regular,BoldFont=*-bold,ItalicFont=*-italic,BoldItalicFont=*-bolditalic]
\setsansfont{texgyreheros}[Extension=.otf,UprightFont=*-regular,BoldFont=*-bold,ItalicFont=*-italic,BoldItalicFont=*-bolditalic]
\setmonofont{texgyrecursor}[Extension=.otf,UprightFont=*-regular,BoldFont=*-bold,ItalicFont=*-italic,BoldItalicFont=*-bolditalic,Scale=MatchLowercase]
\usepackage{amsmath,amssymb}
\usepackage{unicode-math}
\setmathfont{texgyrepagella-math.otf}
\usepackage{xcolor,microtype,enumitem,needspace,fancyhdr}
\usepackage{bookmark}
\definecolor{ink}{HTML}{17344B}
\definecolor{accent}{HTML}{197D80}
\definecolor{muted}{HTML}{596773}
\hypersetup{colorlinks=true,linkcolor=accent,urlcolor=accent,pdftitle={SP-13 - Trace-norm-small
perturbations preserve Hermitian spectral
distributions},pdfauthor={Open Problems in Numerical Linear Algebra}}
\urlstyle{same}
\setlength{\parindent}{0pt}
\setlength{\parskip}{5pt plus 1pt minus 1pt}
\linespread{1.04}
\setlength{\emergencystretch}{3em}
\widowpenalty=10000
\clubpenalty=10000
\displaywidowpenalty=10000
\allowdisplaybreaks[1]
\setlist{nosep,leftmargin=1.5em}
\providecommand{\tightlist}{\setlength{\itemsep}{0pt}\setlength{\parskip}{0pt}}
\providecommand{\pandocbounded}[1]{#1}
\pagestyle{fancy}
\fancyhf{}
\fancyhead[L]{\sffamily\footnotesize\color{muted}OPEN PROBLEMS IN NUMERICAL LINEAR ALGEBRA}
\fancyhead[R]{\sffamily\bfseries\footnotesize\color{accent}SP-13}
\fancyfoot[L]{\parbox[b]{0.9\textwidth}{\sffamily\fontsize{8}{10}\selectfont\color{muted}Editorial ratings; a bounded search cannot certify that a problem remains unsolved.\\Literature check: 2026-09-11}}
\fancyfoot[R]{\sffamily\footnotesize\color{muted}\thepage}
\renewcommand{\headrulewidth}{0.3pt}
\renewcommand{\headrule}{\hbox to\headwidth{\color{accent}\leaders\hrule height \headrulewidth\hfill}}
\makeatletter
\renewcommand{\section}{\Needspace{5\baselineskip}\@startsection{section}{1}{0pt}{12pt plus 2pt minus 2pt}{4pt}{\sffamily\bfseries\large\color{ink}}}
\renewcommand{\subsection}{\Needspace{4\baselineskip}\@startsection{subsection}{2}{0pt}{9pt plus 1pt minus 1pt}{3pt}{\sffamily\bfseries\normalsize\color{ink}}}
\makeatother
\setcounter{secnumdepth}{0}
\begin{document}
{\sffamily\small\color{muted}Eigenvalues and inverse problems\par}
\vspace{3pt}
{\raggedright\hyphenpenalty=10000\exhyphenpenalty=10000\sffamily\bfseries\fontsize{21}{24}\selectfont\color{ink}Trace-norm-small
perturbations preserve Hermitian spectral distributions\par}
\vspace{7pt}
{\sffamily\small\textbf{Difficulty:} challenging\quad\textbf{Importance:} interesting
to the community\par}
{\sffamily\small\bfseries\color{accent}Status: Partially resolved\par}
\vspace{4pt}
{\color{accent}\hrule height 0.8pt}
\vspace{6pt}
\textbf{Topic:} Non-Hermitian perturbations and spectral distributions

\textbf{Rating rationale:} Removing all spectral-norm bounds requires
control of nonnormal spectra under potentially large perturbations. The
result would extend spectral analysis of discretization matrices and
preconditioners with unbounded coefficients.

\subsection{Statement}\label{statement}

For a sequence \(A_n\in\mathbb C^{n\times n}\) and a measurable function
\(f:[0,1]\to\mathbb R\), write \(\{A_n\}\sim_\lambda f\) if

\[
\lim_{n\to\infty}\frac1n\sum_{j=1}^n F(\lambda_j(A_n))
=\int_0^1 F(f(t))\,dt
\qquad\text{for every }F\in C_c(\mathbb C),
\]

where \(C_c(\mathbb C)\) denotes the continuous complex-valued functions
of compact support, and eigenvalues are counted with algebraic
multiplicity. Define the trace norm by \(\|E\|_*:=\sum_j\sigma_j(E)\).

\textbf{Conjecture (Barbarino--Serra-Capizzano).} For every Hermitian
sequence \(H_n=H_n^*\) with \(\{H_n\}\sim_\lambda f\), and every
sequence \(E_n\in\mathbb C^{n\times n}\) satisfying \(\|E_n\|_*/n\to0\),
one has

\[\{H_n+E_n\}\sim_\lambda f.\]

Neither sequence is assumed uniformly bounded in spectral norm;
\(H_n+E_n\) need not be normal. The unit interval is a distributional
normalization of the source's general finite-measure symbol domain.

\subsection{Known cases and numerical
significance}\label{known-cases-and-numerical-significance}

The conclusion is proved if \(\|E_n\|_F=o(\sqrt n)\), and also if
\(\|E_n\|_*=o(n)\) and \(\sup_n\|E_n\|_2<\infty\). Here \(\|\cdot\|_F\)
and \(\|\cdot\|_2\) denote Frobenius and spectral norms. The conjecture
retains the trace-norm assumption while removing the latter uniform
bound. It concerns the limiting empirical distribution, not
individual-eigenvalue matching.

Spectral symbols describe large matrix sequences arising in PDE
discretization and preconditioning. The source's equivalent GLT and
diagonal formulations are grouped with this target rather than assigned
separate IDs.

\newpage
\subsection{References and status
check}\label{references-and-status-check}

\begin{itemize}
\tightlist
\item
  G. Barbarino and S. Serra-Capizzano, \emph{Non-Hermitian perturbations
  of Hermitian matrix-sequences and applications to the spectral
  analysis of the numerical approximation of partial differential
  equations}, Numerical Linear Algebra with Applications 27 (2020),
  e2286. \href{https://doi.org/10.1002/nla.2286}{DOI};
  \href{https://giovannibarbarino.github.io/doc/articles/NHperturbation.pdf}{author
  PDF}. Conjecture 1 in §6, printed p.29; Theorem 1 and Corollary 3 give
  the stated partial results.
\item
  G. Barbarino, \emph{Conjectures on Perturbations of Hermitian
  Sequences},
  \href{https://arxiv.org/abs/1808.05555v1}{arXiv:1808.05555v1},
  Conjecture 1 and Lemma 4.2. This supplies equivalent versions of the
  same target.
\item
  G. Barbarino and C. Garoni, \emph{GLT sequences and normal matrices},
  Electronic Journal of Linear Algebra 41 (2025), 1--20.
  \href{https://journals.uwyo.edu/index.php/ela/article/download/8929/6949/23285}{Theorem
  3.2} requires smaller perturbations when the sum is not normal, so it
  does not settle this conjecture.
\end{itemize}

On 2026-09-11, checked these primary statements, the authors'
publication lists, and targeted title, trace-norm,
spectral-distribution, proof and counterexample searches. No full
resolution was located. The surviving evidence is a bounded literature
check; the 2025 normal-sequence theorem is not presented as a new proof
of openness.
\end{document}
